\begin{tikzpicture}
\node[note] at (0,1.7) {修改历史};
\node[note] at (4.1,1.7) {副本A的区间};
\node[note] at (7.4,1.7) {副本B的区间};
\node[note] at (10.5,1.7) {读取性质};
\foreach \y/\caseName/\bytesA/\bytesB/\caseResult in {
 .5/无并发的成功写/AAAA/AAAA/一致且已定义,
 -1.1/大写入拆分后交错/AABB/AABB/一致但未定义,
 -2.7/修改只到达部分副本/XXXX/AAAA/不一致}
 {\node[note,text width=3cm] at (0,\y) {\caseName};
  \node[cell,minimum width=2.5cm] at (4.1,\y) {\ttfamily\bytesA};
  \node[cell,minimum width=2.5cm] at (7.4,\y) {\ttfamily\bytesB};
  \node[note,text width=2.6cm] at (10.5,\y) {\caseResult};}
\node[note,text width=11.6cm] at (5.1,-4.1) {A、B、X分别代表不同写入的字节。相同次序能使副本一致，\\大写入被拆分时，完整区间可能混合不同客户端的内容。};
\end{tikzpicture}
